\documentclass[12pt]{article}
\usepackage[utf8]{inputenc}
\usepackage[T1]{fontenc}
\usepackage{hyperref}
\title{Kaitag Language Resources for K-12 Curriculum Development}
\author{Compiled by Claude Sonnet 5.5 (claude-sonnet-5-5) for Andr\'as Kornai}
\date{}
\begin{document}
\maketitle

\section{Introduction}
Kaitag (also Kajtak; ISO 639-3 xdq; Glottocode kajt1238) is a Dargin-branch Northeast Caucasian language of Dagestan. Search results report about 30,000 speakers (2020). It is written in Cyrillic and has sometimes been treated as a dialect of Dargwa. The dialects are Lower Kaitag, Upper Kaitag and Shari, which may be a separate language.

\section{Primary Educational Resources}
I found no textbooks or school programmes for Kaitag.

\section{Linguistic and Reference Materials}
\begin{itemize}
\item Magomed Magomedov, a bilingual Kaitag-English dictionary with 5,171 lexemes, based on Gasanova's 2012 dissertation. It is listed on Academia.edu, so it should be checked before use. \url{https://independent.academia.edu/alkaitagi}
\item Wikipedia, ``Kaitag language''. \url{https://en.wikipedia.org/wiki/Kaitag_language}
\item Wikipedia, ``Kaitags'', on the people. \url{https://en.wikipedia.org/wiki/Kaitags}
\item A journal article on a Khaydak-Russian dictionary project for the dialect of the Kaitagsky district. \url{https://en.nbpublish.com/library_read_article.php?id=76193}
\end{itemize}

\section{Religious and Community Resources}
None found with a verified URL.

\section{Digital and Audio Resources}
\begin{itemize}
\item Omniglot, ``Kaitag'', with a writing and orthography overview. \url{https://omniglot.com/writing/kaitag.htm}
\item Global Voices, an article on Dagestan's indigenous languages, for context. \url{https://globalvoices.org/2021/03/28/the-watchtower-on-the-mountain-of-dagestans-indigenous-languages/}
\end{itemize}

\section{Curriculum Development Context}
A dictionary and a few descriptions exist, but I found no standalone grammar or school materials. The language is reported as not taught in schools. A curriculum would need to build primers and a standard orthography on the Cyrillic basis, working with Dagestani linguists and the community.
\end{document}